\chapter{Metodologia}\label{cap:metodologia}

Este capítulo descreve o percurso metodológico adotado para atingir o objetivo geral enunciado no Capítulo~\ref{cap:introducao}, o de desenvolver e avaliar uma abordagem baseada em GraphRAG para a extração de informações clínicas de prontuários eletrônicos de pacientes hepatopatas. Para tornar essa avaliação mensurável, a extração foi formulada como uma tarefa de classificação multirrótulo: dado o texto de um prontuário, o sistema deve indicar quais, dentre oito doenças hepáticas definidas previamente, estão claramente documentadas como presentes no paciente. A pergunta que orienta o trabalho não se limita a verificar se o sistema acerta, mas busca identificar qual configuração de busca do Microsoft GraphRAG apresenta melhor desempenho nesse tipo de extração clínica e por quais razões.

As seções seguem a ordem cronológica de execução do trabalho, sintetizada na Figura~\ref{fig:met-visao-geral}. Inicialmente, apresentam-se os dados e as classes-alvo, seguidos da construção da amostra avaliada e de seus rótulos de referência. Em seguida, descrevem-se a construção do grafo de conhecimento, os modos de busca e os sistemas comparados, o \textit{prompt} de extração, o modelo de linguagem e a infraestrutura de execução. Por fim, detalham-se a conversão das respostas em classes, as métricas e os testes estatísticos, a análise qualitativa de erros e as limitações metodológicas. Os valores de desempenho obtidos são apresentados no capítulo de resultados; neste capítulo constam apenas os números que fundamentam decisões de projeto.

\begin{figure}[!htbp]
\centering
\resizebox{\textwidth}{!}{% Arquitetura do sistema avaliado, em quatro painéis.
% Gramática visual inspirada na figura de visão geral de Xonto-RAG
% (Wang et al., 2026): painéis numerados com fundo pastel, mini-grafos e
% setas largas entre etapas; cores da paleta da tese (estilo-visual.tex).
\begin{tikzpicture}[
  x=1cm, y=1cm, font=\sffamily\footnotesize,
  painel/.style={rounded corners=9pt, line width=0.8pt, draw=black!35},
  titulo/.style={anchor=west, font=\rmfamily\bfseries\normalsize, inner sep=0pt},
  num/.style={circle, draw=black!75, line width=0.8pt, inner sep=1.2pt, minimum size=5mm, font=\rmfamily\bfseries\normalsize},
  caixa/.style={rounded corners=3pt, line width=0.6pt, align=center, inner sep=3pt, fill=white},
  seta/.style={-{Latex[length=1.8mm]}, line width=0.7pt, black!70},
  setao/.style={single arrow, shape border uses incircle, draw=black!45, fill=black!25, minimum height=6.5mm, minimum width=7mm, single arrow head extend=1.3mm, inner sep=0pt},
  no/.style={circle, minimum size=2.6mm, inner sep=0pt, line width=0.6pt},
  na/.style={no, fill=vAzul!45, draw=vAzul!80!black},
  nl/.style={no, fill=vLaranja!50, draw=vLaranja!80!black},
  nv/.style={no, fill=vVerde!45, draw=vVerde!80!black},
  aresta/.style={line width=0.6pt, black!45},
  rotulo/.style={font=\sffamily\scriptsize, text=black!65}
]
\def\PW{7.55}\def\PH{4.55}\def\GX{0.70}\def\GY{0.50}

% ======================================================= fundos dos painéis
\path[painel, fill=vAmbar!14]  (0,0) rectangle (\PW,-\PH);
\path[painel, fill=vAzul!11]   (\PW+\GX,0) rectangle (2*\PW+\GX,-\PH);
\path[painel, fill=vVerde!12]  (0,-\PH-\GY) rectangle (\PW,-2*\PH-\GY);
\path[painel, fill=vRoxo!9]    (\PW+\GX,-\PH-\GY) rectangle (2*\PW+\GX,-2*\PH-\GY);

% setas largas entre painéis
\node[setao] at ({\PW+\GX/2},{-\PH/2}) {};
\node[setao, shape border rotate=225, minimum height=5.5mm] at ({\PW+\GX/2},{-\PH-\GY/2}) {};
\node[setao] at ({\PW+\GX/2},{-1.5*\PH-\GY}) {};

% ============================== 1. Construção da amostra e dos rótulos
\begin{scope}
\node[num] at (0.42,-0.42) {1};
\node[titulo] at (0.82,-0.42) {Construção da amostra};
% funil: faixas cada vez mais estreitas
\foreach \i/\w/\cor/\rot in {
    0/6.7/vNeutro!85/{\faDatabase\ \ Corpus: 3.757 prontuários, 6.405 atendimentos},
    1/5.9/vAmbar!30/{\faFilter\ \ Triagem por expressões regulares},
    2/5.1/vAmbar!45/{\faRobot\ \ Triagem por LLM},
    3/4.3/vAmbar!62/{\faUserMd\ \ Anotação manual (segunda opinião)}}{
  \pgfmathsetmacro{\yt}{-0.85-\i*0.60}
  \pgfmathsetmacro{\yb}{\yt-0.53}
  \pgfmathsetmacro{\wb}{\w-0.8}
  \path[fill=\cor, draw=white, line width=1pt]
    ({\PW/2-\w/2},\yt) -- ({\PW/2+\w/2},\yt) -- ({\PW/2+\wb/2},\yb) -- ({\PW/2-\wb/2},\yb) -- cycle;
  \node[font=\sffamily\scriptsize] at ({\PW/2},{(\yt+\yb)/2}) {\rot};
}
\draw[seta] ({\PW/2},-3.22) -- ({\PW/2},-3.47);
\node[caixa, draw=vAmbar!80!black, text width=6.2cm, font=\sffamily\scriptsize] at ({\PW/2},-3.92)
  {\textbf{537 prontuários} com rótulos de referência\\[1pt]
   \tikz[baseline=-0.6ex]{\foreach \k/\c in {0/vAzul,1/vVerde,2/vLaranja,3/vAmbar,4/vRoxo,5/vVermelho,6/vTeal,7/vCinza}{\fill[\c] (\k*0.32,0) circle (1.1mm);}}
   \ \ oito doenças (multirrótulo)};
\end{scope}

% ======================================================= 2. Indexação (GraphRAG)
\begin{scope}[xshift={(\PW+\GX)*1cm}]
\node[num] at (0.42,-0.42) {2};
\node[titulo] at (0.82,-0.42) {Indexação pelo GraphRAG};
\node[caixa, draw=vAzul!75!black, text width=1.95cm, font=\sffamily\scriptsize] (pt)  at (1.25,-1.15) {Ajuste de\\\textit{prompts}};
\node[caixa, draw=vAzul!75!black, text width=1.95cm, font=\sffamily\scriptsize] (seg) at (3.78,-1.15) {Segmentação\\5.541 unidades};
\node[caixa, draw=vAzul!75!black, text width=1.95cm, font=\sffamily\scriptsize] (ext) at (6.30,-1.15) {Entidades e\\relações (LLM)};
\draw[seta] (pt) -- (seg); \draw[seta] (seg) -- (ext);
% mini-grafo com três comunidades (achatado na vertical)
\begin{scope}[shift={(0.35,-1.90)}, yscale=0.60]
  \draw[vAzul!70, dashed, line width=0.5pt, rounded corners=5pt] (0.05,-0.15) rectangle (1.35,-1.55);
  \draw[vLaranja!80, dashed, line width=0.5pt, rounded corners=5pt] (1.55,-0.15) rectangle (2.95,-1.40);
  \draw[vVerde!80, dashed, line width=0.5pt, rounded corners=5pt] (0.70,-1.75) rectangle (2.45,-2.85);
  \node[na] (a1) at (0.35,-0.45) {}; \node[na] (a2) at (1.05,-0.50) {}; \node[na] (a3) at (0.45,-1.20) {}; \node[na] (a4) at (1.10,-1.25) {};
  \node[nl] (b1) at (1.90,-0.45) {}; \node[nl] (b2) at (2.65,-0.55) {}; \node[nl] (b3) at (2.25,-1.10) {};
  \node[nv] (c1) at (1.00,-2.10) {}; \node[nv] (c2) at (1.75,-2.05) {}; \node[nv] (c3) at (1.35,-2.55) {}; \node[nv] (c4) at (2.15,-2.55) {};
  \draw[aresta] (a1)--(a2) (a1)--(a3) (a2)--(a4) (a3)--(a4) (a2)--(a3) (b1)--(b2) (b2)--(b3) (b1)--(b3)
                (c1)--(c2) (c1)--(c3) (c2)--(c4) (c3)--(c4) (c2)--(c3);
  \draw[aresta, dashed, black!35] (a2)--(b1) (a4)--(c1) (b3)--(c2);
\end{scope}
\node[align=center, rotulo] at (1.85,-4.08) {Grafo: 26.966 entidades,\\88.347 relações};
\draw[seta] (ext.south) -- ++(0,-0.20) -| (2.45,-1.97);
% hierarquia de Leiden
\begin{scope}[shift={(4.75,-1.95)}]
  \node[rotulo, anchor=east] at (0.05,0) {nível 0};
  \node[rotulo, anchor=east] at (0.05,-0.58) {nível 1};
  \node[rotulo, anchor=east] at (0.05,-1.16) {nível 2};
  \node[no, fill=vNeutro, draw=black!55, minimum size=3.4mm] (r) at (1.35,0) {};
  \node[na, minimum size=3mm] (h1) at (0.65,-0.58) {}; \node[nl, minimum size=3mm] (h2) at (1.35,-0.58) {}; \node[nv, minimum size=3mm] (h3) at (2.05,-0.58) {};
  \foreach \k/\x/\s in {1/0.40/na, 2/0.90/na, 3/1.20/nl, 4/1.55/nl, 5/1.85/nv, 6/2.30/nv}{\node[\s, minimum size=2.2mm] (t\k) at (\x,-1.16) {};}
  \draw[aresta] (r)--(h1) (r)--(h2) (r)--(h3) (h1)--(t1) (h1)--(t2) (h2)--(t3) (h2)--(t4) (h3)--(t5) (h3)--(t6);
  \node[rotulo] at (1.35,-1.45) {$\vdots$ até o nível 6};
\end{scope}
\node[align=center, rotulo] at (5.80,-4.08) {Comunidades de Leiden\\e resumos de comunidade};
\draw[seta] (3.32,-2.75) -- (3.82,-2.75);
\end{scope}

% ======================================================= 3. Consulta (S1 a S8)
\begin{scope}[yshift={-(\PH+\GY)*1cm}]
\node[num] at (0.42,-0.42) {3};
\node[titulo] at (0.82,-0.42) {Consulta: sistemas S1 a S8};
\node[caixa, draw=vVerde!75!black, text width=1.95cm, font=\sffamily\scriptsize] (pr) at (1.22,-1.55) {\faFileMedical\ \ Prontuário\\{\tiny até 8.000 caracteres}};
\node[caixa, draw=vVerde!75!black, text width=1.95cm, font=\sffamily\scriptsize] (pp) at (1.22,-2.60) {\faAlignLeft\ \ \textit{Prompt}\\{\tiny de extração}};
% grade modo x nível
\begin{scope}[shift={(2.62,-0.90)}]
  \node[rotulo] at (2.175,0.02) {nível de comunidade};
  \node[font=\sffamily\scriptsize\bfseries] at (1.75,-0.28) {1};
  \node[font=\sffamily\scriptsize\bfseries] at (2.60,-0.28) {2};
  \foreach \j/\m/\c/\a/\b/\op in {1/basic/modoBasic/S1/S2/1, 2/local/modoLocal/S3/S4/1, 3/global/modoGlobal/S5/S6/0.4, 4/drift/modoDrift/S7/S8/1}{
    \pgfmathsetmacro{\y}{-0.28-\j*0.47}
    \node[anchor=west, inner sep=0pt] at (0,\y) {\chip{\c}\textit{\m}};
    \node[caixa, draw=\c!80!black, fill=\c!20, minimum width=8mm, inner sep=2pt, opacity=\op, text opacity=\op] at (1.75,\y) {\a};
    \node[caixa, draw=\c!80!black, fill=\c!20, minimum width=8mm, inner sep=2pt, opacity=\op, text opacity=\op] at (2.60,\y) {\b};
  }
\end{scope}
\draw[line width=0.7pt, black!70] (pr.east) -| (2.40,-2.08);
\draw[line width=0.7pt, black!70] (pp.east) -| (2.40,-2.08);
\draw[seta] (2.40,-2.08) -- (2.60,-2.08);
\node[caixa, draw=vVerde!75!black, text width=1.0cm] (llm) at (6.60,-2.03) {{\Large\faRobot}\\[1pt] LLM};
\node[font=\ttfamily\tiny, text=black!70, anchor=north] at (llm.south) {gpt-oss-120b};
\draw[seta] (5.64,-2.03) -- (llm.west);
\node[fill=vVermelho!75, text=white, rounded corners=2pt, align=left, inner sep=3.5pt, font=\sffamily\scriptsize] at ({\PW/2},-3.97)
  {\textbf{Pergunta:} quais das oito doenças hepáticas estão\\documentadas como presentes neste prontuário?};
\end{scope}

% ======================================================= 4. Extração e avaliação
\begin{scope}[xshift={(\PW+\GX)*1cm}, yshift={-(\PH+\GY)*1cm}]
\node[num] at (0.42,-0.42) {4};
\node[titulo] at (0.82,-0.42) {Extração e avaliação};
\node[caixa, draw=black!45, text width=3.25cm, align=left, font=\ttfamily\tiny] (resp) at (1.95,-1.62)
  {{\sffamily\scriptsize\bfseries Resposta do modelo}\\[2pt] <RESPOSTA>\\ Cirrose hepática: ``\ldots''\\ Hipertensão portal: ``\ldots''\\ </RESPOSTA>};
\node[caixa, draw=vRoxo!75!black, text width=2.75cm, font=\sffamily\scriptsize] (parser) at (5.75,-1.62)
  {\faCogs\ \textbf{\textit{Parser}}\\determinístico\\[1pt]{\tiny\color{black!65}nomes canônicos e negação}};
\draw[seta] (resp) -- (parser);
\node[caixa, draw=vRoxo!75!black, text width=6.75cm, align=left, font=\sffamily\scriptsize] (aval) at ({\PW/2},-3.45)
  {\textbf{Doenças previstas} $\times$ \textbf{rótulos de referência}\\[1pt]
   \faChartBar\ \ métricas multirrótulo (EM, P, R, F1)\\
   \faFlask\ \ hipóteses H1 a H4 sobre o modo \textit{drift}\\
   \faSearch\ \ análise qualitativa de erros e do grafo};
\draw[seta] (parser.south) -- (parser.south |- aval.north);
\end{scope}
\end{tikzpicture}
}
\caption[Arquitetura do sistema avaliado]{Arquitetura do sistema avaliado: (1) construção da amostra (Seção~\ref{sec:met-amostra}); (2) indexação pelo GraphRAG (Seção~\ref{sec:met-grafo}); (3) consulta pelos sistemas S1 a S8, com S5 e S6 esmaecidos por não terem sido concluídos (Seção~\ref{sec:met-sistemas}); e (4) extração e avaliação.}
\label{fig:met-visao-geral}
\par\smallskip{\footnotesize Fonte: elaborado pelo autor.}
\end{figure}

\section{Dados}\label{sec:met-dados}

Os dados utilizados são prontuários eletrônicos reais do ambulatório de hepatologia do Hospital Universitário da Universidade Federal de Alagoas (HU-UFAL), em Maceió, Alagoas.
O corpus reúne 3.757 prontuários únicos, que totalizam 6.405 registros de atendimento. Cada prontuário corresponde à concatenação, em ordem cronológica, de todas as evoluções clínicas de um mesmo paciente, com os atendimentos individuais delimitados pelo marcador textual \texttt{-{}-{}- Atendimento -{}-{}-}. O conteúdo é predominantemente texto livre e exibe as características discutidas na Seção~\ref{sec:prontuarios}, como abreviações não padronizadas, estilo telegráfico e menções frequentes a hipóteses negadas ou não confirmadas.

Quanto aos aspectos éticos, os prontuários foram anonimizados antes de qualquer uso neste trabalho.
% TODO: descrever o procedimento de anonimização (ferramenta, categorias de dados removidas, verificação do resultado).
Essa etapa é particularmente relevante porque, nas consultas da rodada final, o texto anonimizado dos prontuários foi enviado a um provedor externo de inferência, acessado por meio de uma interface de programação de aplicações (API), conforme a Seção~\ref{sec:met-modelo}, e não processado exclusivamente em ambiente local. A anonimização prévia constitui, portanto, a principal salvaguarda adotada em relação ao tratamento de dados pessoais sensíveis disciplinado pela Lei Geral de Proteção de Dados Pessoais (LGPD).
% TODO: citar o texto da LGPD (não há entrada no .bib).
% TODO: informar parecer do CEP/CAAE.

\section{Definição das classes-alvo}\label{sec:met-classes}

As oito classes avaliadas são cirrose hepática, esquistossomose, hepatite B, hepatite C, esteatose hepática, hipertensão portal, hemorragia digestiva alta (HDA) e peritonite bacteriana espontânea (PBE). Essas condições não foram selecionadas por critério estatístico nem descobertas de forma exploratória no texto: foram definidas \textit{a priori} por uma profissional de hepatologia, com base em sua relevância clínica no contexto do ambulatório de origem dos dados.

\label{sec:estrutura-causal}As oito condições não são independentes entre si, e sua estrutura causal condiciona a formulação da tarefa. A esteatose hepática, a hepatite B e a hepatite C são causas frequentes de dano hepático crônico que pode evoluir para cirrose. A cirrose e a esquistossomose hepatoesplênica são causas de hipertensão portal, respectivamente nas formas cirrótica e não cirrótica. A hemorragia digestiva alta e a peritonite bacteriana espontânea são descompensações da doença hepática de base, associadas à hipertensão portal e à cirrose avançada. Um mesmo paciente pode, portanto, apresentar várias dessas condições encadeadas, o que motiva a formulação da tarefa como classificação multirrótulo: a cada prontuário associa-se o subconjunto das oito condições documentadas como presentes, em vez de uma única categoria. Ainda assim, a estrutura causal exige cuidado na distinção entre doença de base e complicação, de modo que a fronteira entre classes permanece um ponto sensível da avaliação, discutido na definição do critério de presença clara na Seção~\ref{sec:met-prompt}.

A lista fixa de classes é o ponto de partida de todo o pipeline. Tanto a triagem de prontuários candidatos, descrita na Seção~\ref{sec:met-amostra}, quanto o \textit{prompt} de extração, apresentado na Seção~\ref{sec:met-prompt}, foram construídos em torno exatamente dessas oito categorias. Definidas as classes, o passo seguinte consistiu em selecionar os prontuários a avaliar e estabelecer, para cada um deles, o conjunto correto de doenças.

\section{Construção da amostra e dos rótulos de referência}\label{sec:met-amostra}

Os experimentos foram conduzidos sobre um conjunto de 537 prontuários extraído do corpus de 3.757. Esse conjunto não constitui amostra aleatória: foi construído por um funil de três etapas, representado no painel 1 da Figura~\ref{fig:met-visao-geral}. A escolha decorre do fato de que uma amostra aleatória simples teria prevalência baixa demais de cada uma das oito doenças, com número de casos positivos por classe insuficiente para uma avaliação multirrótulo. O funil contorna esse problema ao enriquecer a amostra com prontuários que plausivelmente mencionam as condições de interesse.

A primeira etapa consistiu em uma triagem por expressões regulares sobre o texto bruto de todo o corpus, em busca de menções às oito doenças-alvo. Um teste manual preliminar mostrou que uma busca ingênua pelos termos produzia falsos positivos sistemáticos, o que motivou quatro correções:
\begin{enumerate}
    \item desconsiderar a colisão entre a sigla HDA da doença-alvo e o cabeçalho de seção ``H.D.A.'' (história da doença atual), que respondia sozinha por cerca de 73\% dos falsos positivos de HDA antes da correção;
    \item descartar menções explicitamente negadas, marcadas por termos como ``nega'', ``ausente'' e ``descartado'';
    \item descartar menções marcadas por incerteza (\textit{hedges}), como ``?'', ``a esclarecer'', ``suspeita de'' e ``epidemiologia positiva'', que indicam risco, e não achado positivo;
    \item descartar menções em contexto histórico ou já resolvido, como ``passado de'' e ``história de'', exceto quando a condição é o motivo do atendimento atual.
\end{enumerate}
A saída dessa etapa foi uma lista de prontuários candidatos, isto é, que mencionam ao menos uma das oito doenças de forma textualmente plausível.

A segunda etapa submeteu os candidatos a uma triagem por LLM. Para cada prontuário, realizou-se uma chamada direta ao modelo, sem indexação nem recuperação, com uma pergunta curta que solicitava confirmação binária (sim ou não) para cada doença candidata. O objetivo foi reduzir o volume de falsos positivos antes da revisão humana, a etapa mais onerosa do funil. As chamadas foram executadas em paralelo e de forma retomável, o que permitiu reiniciar o processamento a partir do ponto de interrupção.
% TODO: informar o modelo de linguagem usado na triagem por LLM.

A terceira etapa foi a anotação manual, denominada no projeto ``segunda opinião'': uma revisão humana dos candidatos que passaram pelas etapas anteriores, que produziu o rótulo definitivo de cada prontuário. Esse conjunto constitui o \textit{ground truth} do trabalho, isto é, o conjunto de rótulos considerados corretos contra os quais as saídas do sistema são comparadas, e é referido daqui em diante como rótulos de referência.
% TODO: informar quantos anotadores participaram, sua formação e se houve medida de concordância entre anotadores (por exemplo, kappa de Cohen). A Seção 2.7 afirma que o trabalho adota ground truth "anotado por especialista"; confirmar.

Os números disponíveis confirmam que o funil molda a composição da amostra final. Em uma rodada específica de expansão da amostra, foram registrados 396 prontuários candidatos. Dos 68 prontuários que terminaram rotulados sem nenhuma das oito doenças, 43 (63\%) pertencem a esse grupo de candidatos, ou seja, foram sinalizados pela triagem textual, mas a revisão por LLM ou manual concluiu que a menção não correspondia a diagnóstico presente, seja por negação, por incerteza ou por contexto histórico não capturados pelas expressões regulares. O grupo sem doença, portanto, não é um grupo de controle aleatório: resulta, em parte substancial, da própria filtragem de falsos positivos.

A Tabela~\ref{tab:met-distribuicao} apresenta a distribuição final das classes. Por se tratar de classificação multirrótulo, os percentuais não somam 100\%, e cada paciente apresenta, em média, 1,49 doenças confirmadas. Duas ressalvas decorrem dessa distribuição. A primeira refere-se à PBE, com apenas 2 casos positivos em toda a amostra: qualquer métrica calculada sobre essa classe deve ser lida como ruído estatístico, e não como indicação de desempenho, pois uma única previsão a mais ou a menos altera o valor por completo. A segunda refere-se à prevalência observada, que é artefato do critério de seleção enriquecido e não deve ser extrapolada como estimativa epidemiológica da população atendida no ambulatório.

\begin{table}[!htbp]
\centering
\caption[Distribuição das classes na amostra de 537 prontuários]{Distribuição das classes na amostra de 537 prontuários. Como a classificação é multirrótulo, os percentuais não somam 100\%.}
\label{tab:met-distribuicao}
\begin{tabular}{lrl}
\toprule
Doença & Pacientes & \% da amostra \\
\midrule
Hipertensão portal & 201 & \barcell{vTeal}{0.9350}{37,4} \\
Esteatose hepática & 156 & \barcell{vTeal}{0.7275}{29,1} \\
Cirrose hepática & 152 & \barcell{vTeal}{0.7075}{28,3} \\
Esquistossomose & 95 & \barcell{vTeal}{0.4425}{17,7} \\
Hepatite B & 90 & \barcell{vTeal}{0.4200}{16,8} \\
Hepatite C & 67 & \barcell{vTeal}{0.3125}{12,5} \\
Hemorragia digestiva alta & 36 & \barcell{vTeal}{0.1675}{6,7} \\
Peritonite bacteriana espontânea & 2 & \barcell{vTeal}{0.0100}{0,4} \\
\midrule
Nenhuma das oito doenças & 68 & \barcell{vCinza}{0.3175}{12,7} \\
\bottomrule
\end{tabular}
\par\smallskip{\footnotesize Fonte: elaborado pelo autor.}
\end{table}

\section{Construção do grafo de conhecimento}\label{sec:met-grafo}

Paralelamente à construção da amostra, o corpus foi indexado pelo GraphRAG. O índice foi construído uma única vez, antes de qualquer consulta, segundo o pipeline de indexação descrito na Seção~\ref{sec:graphrag} e representado no painel 2 da Figura~\ref{fig:met-visao-geral}. Esta seção não retoma o funcionamento geral de cada etapa, mas registra como ela foi aplicada ao corpus e quais estruturas resultaram.
% TODO: confirmar se o índice cobre o corpus completo de 3.757 prontuários ou apenas a amostra de 537.
% TODO: informar o modelo de linguagem e o modelo de embeddings usados na indexação, o tamanho das unidades de texto e a sobreposição entre elas.

Antes da indexação, executou-se o procedimento automático de ajuste de \textit{prompts} (\textit{prompt tuning}) oferecido pelo próprio GraphRAG, que gera \textit{prompts} de extração adaptados ao domínio a partir de uma amostra do corpus. A diferença mais relevante entre as versões genérica e ajustada está no \textit{prompt} de extração de entidades e relações. A versão genérica exige uma lista fixa de tipos de entidade, tipicamente categorias amplas de propósito geral, como pessoa, organização e local; a versão ajustada não fixa os tipos e solicita ao modelo que sugira categorias, tão gerais quanto possível, a partir do próprio conteúdo.

Em decorrência desse ajuste, os tipos de entidade observados no grafo final são específicos do domínio clínico, como \texttt{CONDICAO\_CLINICA}, \texttt{EXAME}, \texttt{EXAME\_LABORATORIAL}, \texttt{MEDICAMENTO}, \texttt{PROCEDIMENTO}, \texttt{SINAL\_CLINICO}, \texttt{ÓRGÃO}, \texttt{HÁBITO} e \texttt{PROTOCOLO}, em vez de categorias genéricas sem relação com o domínio. Como todos os modos de busca dependem, em graus diferentes, da qualidade do grafo, esse ajuste é pré-requisito comum a todos os sistemas avaliados.

O cuidado com essa etapa tem respaldo na literatura. \citet{ye2026} relatam que o pipeline padrão do GraphRAG, aplicado a prontuários eletrônicos, produz granularidade excessiva, que degrada significativamente a recuperação: o \textit{prompt} original identificava cada parâmetro laboratorial como entidade independente, o que provocava uma proliferação acentuada de entidades. Os autores contornaram o problema com \textit{prompts} customizados para o domínio, que incluem a exigência de que cada registro contenha exatamente uma entidade do tipo paciente, à qual se vinculam diagnósticos, medicamentos e exames. Neste trabalho, o ajuste empregado foi o procedimento automático do próprio GraphRAG, sem regras explícitas dessa natureza, o que tem consequências para a estrutura do grafo discutidas na Seção~\ref{sec:res-grafos}.

A indexação propriamente dita compreendeu quatro etapas. Na segmentação (\textit{chunking}), o corpus foi dividido em 5.541 unidades de texto (\textit{text units}), unidade mínima à qual se vinculam tanto as entidades extraídas quanto os documentos de origem. Na extração de entidades e relações, um LLM leu cada unidade de texto e identificou entidades, como diagnósticos, exames, medicamentos, procedimentos e sinais clínicos, bem como as relações entre elas, o que resultou em um grafo com 26.966 entidades e 88.347 relações.

A detecção de comunidades empregou a implementação hierárquica do algoritmo de Leiden \citep{traag2019} disponível na biblioteca graspologic, que particiona o grafo em grupos de nós densamente interconectados, maximizando a modularidade. O caráter recursivo do algoritmo, discutido na Seção~\ref{sec:graphrag}, faz com que as comunidades de um nível sejam subdivididas no nível seguinte, o que gerou uma hierarquia de 7 níveis, numerados de 0, o mais amplo, a 6, o mais granular. Cabe registrar que essa partição é puramente topológica: depende apenas de quais entidades estão conectadas e com que peso, sem nenhuma etapa semântica no agrupamento em si.

Por fim, na geração de resumos de comunidade (\textit{community reports}), um LLM produziu, para cada comunidade, um resumo textual a partir das entidades, das relações e dos trechos de texto associados a ela. Esses resumos constituem o material consultado pelos modos \textit{global} e \textit{drift}. O índice resultante é comum a todos os sistemas; o que os distingue é a forma como cada modo de busca o percorre, tema da seção seguinte.

\section{Modos de busca e sistemas avaliados}\label{sec:met-sistemas}

Com o índice construído, cada consulta, composta pelo texto de um prontuário e pelo \textit{prompt} de extração, pode ser respondida por quatro modos de busca, cujo funcionamento geral foi apresentado na Seção~\ref{sec:graphrag} e na Figura~\ref{fig:modos-busca}. No modo \textit{basic}, não há uso estruturado do grafo: como o texto do prontuário já está embutido na própria consulta, o modo opera essencialmente como uma chamada direta ao modelo de linguagem, sem recuperação adicional a partir do índice.
% TODO: conferir com a Seção 2.5, que descreve o basic como recuperação de unidades de texto por similaridade vetorial; se houve recuperação de unidades de texto nesta configuração, ajustar o texto.
O modo \textit{local} recupera contexto centrado nas entidades relevantes para a consulta, com suas relações diretas e as unidades de texto de origem. Seu desempenho depende diretamente do que foi extraído na indexação, pois uma entidade que não foi criada nessa etapa não pode ser recuperada depois.

O modo \textit{global} percorre os resumos de comunidade, seja de um nível fixo da hierarquia, seja, quando a seleção dinâmica de comunidades (\textit{dynamic community selection}) está ativa, navegando a hierarquia de forma adaptativa a partir do nível 0 e avaliando a relevância de cada comunidade antes de decidir se desce mais um nível. O modo \textit{drift} combina as orientações local e global de forma iterativa: realiza uma varredura inicial (\textit{primer}), gera perguntas de acompanhamento (\textit{follow-ups}) com base no que já foi encontrado, configuradas para até 6 por consulta, e executa buscas locais específicas para cada uma delas, refinando a resposta em múltiplas rodadas.

\subsection{Desenho experimental}\label{sec:met-desenho}

O cruzamento dos quatro modos de busca com dois níveis de comunidade, 1 e 2, produz os oito sistemas apresentados na Tabela~\ref{tab:met-sistemas}. O desenho prevê que todos sejam avaliados lado a lado sobre os mesmos 537 prontuários, com o mesmo \textit{prompt} de extração e o mesmo modelo de linguagem, de forma que o modo de busca seja a principal variável isolada na comparação, e o nível de comunidade, a variável secundária. Os sistemas S5 e S6, correspondentes ao modo \textit{global}, integram o desenho, mas não foram concluídos de forma confiável com o modelo adotado, pelos motivos expostos na Seção~\ref{sec:met-limitacoes}; a comparação efetiva restringe-se, por isso, aos seis sistemas restantes.
% TODO: esclarecer se S5 e S6 usaram seleção dinâmica de comunidades e como o nível 1 ou 2 se aplica nesse caso.

\begin{table}[!htbp]
\centering
\caption[Sistemas avaliados, por modo de busca e nível de comunidade]{Sistemas avaliados, definidos pelo cruzamento entre modo de busca e nível de comunidade.}
\label{tab:met-sistemas}
\begin{tabular}{lcc}
\toprule
Modo de busca & Nível 1 & Nível 2 \\
\midrule
\chip{modoBasic}\textit{basic}  & S1 & S2 \\
\chip{modoLocal}\textit{local}  & S3 & S4 \\
\chip{modoGlobal}\textit{global} & {\color{black!45}S5} & {\color{black!45}S6} \\
\chip{modoDrift}\textit{drift}  & S7 & S8 \\
\bottomrule
\end{tabular}
\par\smallskip{\footnotesize Fonte: elaborado pelo autor. Em cinza, os sistemas não concluídos de forma confiável (Seção~\ref{sec:met-limitacoes}).}
\end{table}

Uma consequência do desenho merece registro. Como o modo \textit{basic} não utiliza a estrutura do grafo, o nível de comunidade não interfere na montagem do seu contexto, de modo que S1 e S2 correspondem, na prática, a duas execuções independentes da mesma configuração. Esse par oferece, assim, uma referência da variação entre execuções atribuível ao próprio modelo, útil para interpretar as diferenças entre níveis observadas nos demais modos.
% TODO: confirmar que, na configuração usada, o parâmetro de nível de comunidade não altera o comportamento do modo basic.

\subsection{Escolha dos níveis de comunidade}\label{sec:met-niveis}

A hierarquia produzida pelo algoritmo de Leiden possui 7 níveis, e a Tabela~\ref{tab:met-niveis} mostra como o número e o tamanho das comunidades variam ao longo dela. O nível 0 foi descartado por conter apenas 46 comunidades, com tamanho médio de 543,6 entidades. A distribuição é fortemente assimétrica, com mediana de apenas 5 entidades, o que indica que poucas comunidades muito grandes concentram a maior parte do grafo. Utilizá-lo agregaria, em uma mesma comunidade, informação de muitos pacientes e de temas heterogêneos, o que é incompatível com uma tarefa que exige granularidade por paciente. Esse perfil é característico das hierarquias produzidas pelo algoritmo de Leiden em grafos de conhecimento, que tendem a ser dominadas por poucas comunidades grandes e cujos tamanhos variam amplamente entre níveis, o que torna irregular a qualidade dos resumos gerados a partir delas \citep{hossain2026}. Os mesmos autores defendem comunidades de tamanho limitado e que preservem a conectividade interna como as mais adequadas à recuperação e à sumarização, critério com o qual é compatível a escolha descrita a seguir.

\begin{table}[!htbp]
\centering
\caption[Número e tamanho das comunidades por nível da hierarquia]{Número e tamanho das comunidades, em entidades, em cada nível da hierarquia do índice. Os níveis 1 e 2 foram os adotados nos sistemas avaliados.}
\label{tab:met-niveis}
\begin{tabular}{crrr}
\toprule
Nível & Comunidades & Tamanho médio & Tamanho mediano \\
\midrule
0 & 46    & 543,6 & 5 \\
\rowcolor{linhaDestaque}1 & 501   & 49,7  & 5 \\
\rowcolor{linhaDestaque}2 & 1.475 & 16,1  & 6 \\
3 & 1.959 & 9,7   & 4 \\
4 & 591   & 15,9  & 3 \\
5 & 60    & 85,1  & 6 \\
6 & 7     & 577,3 & 2 \\
\bottomrule
\end{tabular}
\par\smallskip{\footnotesize Fonte: elaborado pelo autor a partir do índice construído. Em destaque, os níveis adotados.}
\end{table}

Os níveis mais profundos foram descartados por razões opostas. O nível 3 apresenta o maior número de comunidades da hierarquia (1.959), com mediana de 4 entidades, o que amplia o custo computacional dos modos que percorrem comunidades inteiras, como o \textit{global}, e aumenta o risco de separar entidades clinicamente relacionadas em comunidades distintas. Nos níveis 4 a 6, o número de comunidades volta a diminuir, enquanto a relação entre tamanho médio e mediano se torna irregular, sem ganho evidente de especificidade. Os níveis 1 e 2, com 501 e 1.475 comunidades e medianas de 5 e 6 entidades, ocupam a faixa intermediária da hierarquia: são finos o suficiente para separar temas clínicos distintos, sem a fragmentação e o custo dos níveis mais profundos, em consonância com o compromisso entre abrangência e especificidade discutido na Seção~\ref{sec:graphrag}.

Avaliar os dois níveis, em vez de escolher um único \textit{a priori}, permite verificar empiricamente se a escolha do nível afeta o desempenho. Essa verificação foi feita por meio do teste estatístico descrito a seguir.

\subsection{Comparação entre níveis: teste de McNemar exato}\label{sec:met-mcnemar}

Para verificar se a escolha entre os níveis 1 e 2 produz diferença estatisticamente significativa, aplicou-se, para cada modo de busca, o teste de McNemar exato \citep{mcnemar1947} sobre o acerto exato por paciente (\textit{exact match}, definido na Seção~\ref{sec:met-metricas}).
O teste é pareado porque os mesmos pacientes são avaliados nos dois níveis de um mesmo modo. Para cada paciente, registra-se se o conjunto de doenças previsto coincide integralmente com o de referência em cada nível, e apenas os pares discordantes entram na estatística: $b$ denota o número de pacientes acertados somente no nível 1, e $c$, o número de pacientes acertados somente no nível 2.

Sob a hipótese nula de que os dois níveis têm a mesma probabilidade de acerto, cada par discordante tem igual chance de favorecer qualquer um deles, de modo que $b$ segue distribuição binomial com $n = b + c$ tentativas e probabilidade $1/2$. A versão exata do teste calcula o p-valor bilateral diretamente dessa distribuição, conforme a Equação~\eqref{eq:mcnemar}, sem recorrer à aproximação pela distribuição qui-quadrado, o que a torna adequada mesmo quando o número de discordâncias é pequeno. Adotou-se nível de significância de 5\%.
% TODO: no modo drift o teste foi aplicado a 531 pacientes, e não a 537; explicar a diferença (por exemplo, pacientes sem resposta válida em um dos níveis), se confirmado.
\begin{equation}
p = \min\left(1,\; 2\sum_{k=0}^{\min(b,c)} \binom{n}{k}\left(\frac{1}{2}\right)^{n}\right)
\label{eq:mcnemar}
\end{equation}

\section{\texorpdfstring{\textit{Prompt}}{Prompt} de extração}\label{sec:met-prompt}

O \textit{prompt} enviado em cada consulta passou por três versões ao longo do trabalho, cada uma voltada a corrigir um problema concreto identificado na análise de erros das versões anteriores. As três compartilham uma estrutura comum. Todas solicitam explicitamente que o modelo indique quais das oito doenças estão claramente documentadas como presentes, entendendo-se por isso diagnóstico explícito, exame confirmatório ou estadiamento formal. Todas truncam o texto do prontuário em 8.000 caracteres e exigem que a resposta seja delimitada pelas marcas \texttt{<RESPOSTA>} e \texttt{</RESPOSTA>}, com uma doença por linha no formato ``nome da doença: trecho que confirma''.

Integram também a estrutura comum duas regras de interpretação. A regra de negação estabelece que resultados registrados como ``NR'' ou ``negativo'' indicam ausência, e não presença, da condição. A regra temporal, específica para HDA e PBE, determina que essas condições só sejam consideradas quando ativas ou recentes, e não quando se tratar de episódio antigo já resolvido.

A Tabela~\ref{tab:met-prompts} resume a evolução das versões. A versão 2 acrescentou a exigência de que o nome da doença na resposta seja exatamente um dos oito nomes canônicos da lista, sem variação nem sinônimo, mudança motivada pela dificuldade do \textit{parser} (Seção~\ref{sec:met-parser}) em reconhecer as variações livres de nomenclatura produzidas com a primeira versão.

\begin{table}[!htbp]
\centering
\caption[Evolução do \textit{prompt} de extração]{Evolução do \textit{prompt} de extração e problema que motivou cada versão.}
\label{tab:met-prompts}
\begin{tabular}{cp{7.0cm}p{5.6cm}}
\toprule
Versão & Alteração introduzida & Problema que a motivou \\
\midrule
1 & \raggedright Estrutura comum: critério de presença clara, truncamento em 8.000 caracteres, resposta delimitada, regra de negação e regra temporal para HDA e PBE & \raggedright Versão inicial \tabularnewline
2 & \raggedright Nome da doença restrito aos oito nomes canônicos, sem sinônimos & \raggedright Variações livres de nomenclatura não reconhecidas pelo \textit{parser} \tabularnewline
3 & \raggedright Bloco com três confusões recorrentes (HDA e sangramento baixo; hipertensão portal e arterial; escores não escritos e achados sem relação hepática) e reforço contra a listagem por risco ou epidemiologia & \raggedright Erros recorrentes identificados na análise qualitativa \tabularnewline
\bottomrule
\end{tabular}
\par\smallskip{\footnotesize Fonte: elaborado pelo autor.}
\end{table}

A versão 3 acrescentou um bloco explícito com três confusões recorrentes identificadas na análise qualitativa dos erros. A primeira é a distinção anatômica entre hemorragia digestiva alta e sangramento baixo ou retal, como o decorrente de hemorroida ou fissura anal. A segunda é a distinção entre hipertensão portal e hipertensão arterial, termos que compartilham apenas a palavra ``hipertensão''. A terceira é a proibição de citar ou inferir escores de estadiamento, como Child-Pugh, MELD e FIB-4, que não estejam literalmente escritos no texto, bem como de usar como evidência achados de órgãos ou sistemas sem relação hepática direta, salvo quando o texto estabelece essa relação explicitamente. A versão 3 reforçou, ainda, a proibição de listar uma doença apenas por menção de risco, probabilidade ou epidemiologia, sem diagnóstico confirmado.

A versão 3 é a versão final do \textit{prompt} e foi empregada nos sistemas S1 a S4 na rodada definitiva. Uma vez fixado o \textit{prompt}, as consultas passam a depender do modelo de linguagem que as responde e da infraestrutura que as executa, descritos a seguir.
% TODO: informar qual versão do prompt foi usada em S7 e S8 (drift) na rodada final; a fonte registra apenas que a versão 3 foi "testada isoladamente no drift".

\section{Modelo de linguagem e infraestrutura de execução}\label{sec:met-modelo}

A rodada que produziu os resultados reportados no Capítulo~\ref{cap:resultados} -- os oito sistemas sobre os 537 prontuários -- utilizou o modelo \texttt{inclusionai/ling-2.6-flash}, acessado pelo serviço de roteamento OpenRouter. Trata-se de um modelo instrucional, sem etapa de raciocínio expandido, o que mantém a latência baixa e previsível: a duração média observada no modo \textit{drift} (o mais custoso, por encadear múltiplas chamadas) foi de 214 segundos por consulta, com mediana de 190 segundos.
% TODO: informar temperatura e demais parâmetros de decodificação, se fixados, e o provedor/custo por token do ling-2.6-flash na rodada final.

Houve, posteriormente, uma tentativa de migrar todos os sistemas para \texttt{inclusionai/ling-3.0-flash}, um modelo de raciocínio híbrido (\textit{hybrid reasoning}) da mesma família, com o objetivo de uniformizar o modelo entre todas as comparações. Essa tentativa, no entanto, não foi concluída: cobriu apenas os modos \textit{basic} e \textit{local} (S1 a S4), e os modos \textit{global} e \textit{drift} não chegaram a produzir um conjunto de resultados confiável nessa migração, pelos motivos discutidos na Seção~\ref{sec:met-limitacoes}. Os resultados do Capítulo~\ref{cap:resultados}, portanto, referem-se integralmente à rodada em \texttt{ling-2.6-flash}, e não à tentativa de migração para \texttt{ling-3.0-flash}.
% TODO: revisar esta seção à luz da tentativa de migração: decidir se os
% resultados parciais em ling-3.0-flash (S1-S4) devem ser citados no
% capítulo de resultados como comparação complementar, e registrar
% explicitamente que a tentativa de validação do prompt de HDA (Seção
% \ref{sec:res-validacao-hda}) foi conduzida sobre ling-3.0-flash, não sobre
% o modelo da rodada final -- a comparação ali deve ser lida como um
% experimento novo num modelo diferente e mais lento, não como réplica
% exata da condição de S7/S8.

O modelo \texttt{ling-3.0-flash} é um modelo de raciocínio híbrido, que gera \textit{tokens} de raciocínio não exibidos antes de produzir a resposta final, acessado pelo mesmo serviço de roteamento, com inferência realizada pelo provedor Novita ao custo de US\$~0,021 por milhão de \textit{tokens} de entrada e US\$~0,063 por milhão de \textit{tokens} de saída. Esse comportamento aumenta de forma substancial o custo e a latência e torna esta última muito variável: uma mesma consulta pode levar de segundos a horas, de forma pouco previsível -- mesmo no modo \textit{basic} (S1), o mais simples, a duração média observada na tentativa de migração subiu para 566 segundos, cerca de 25 vezes a média equivalente em \texttt{ling-2.6-flash} no mesmo modo. O esforço de raciocínio foi configurado no nível baixo. A execução com esse modelo e com esse provedor exigiu, contudo, três ajustes no código do GraphRAG, descritos nos parágrafos seguintes.

O primeiro ajuste refere-se ao formato de resposta estruturada. O GraphRAG original força respostas no formato de objeto JSON (\textit{JavaScript Object Notation}) nas etapas de \textit{map} do modo \textit{global} e de \textit{primer} do modo \textit{drift}, configuração recusada por provedores que aceitam apenas respostas validadas por esquema JSON. Uma primeira correção removeu essa exigência de todas as chamadas ao LLM, mas afetou também a etapa de avaliação de relevância de comunidades, usada pela seleção dinâmica, e provocou centenas de respostas malformadas por consulta. A correção final restringiu a remoção, por meio da inspeção da pilha de chamadas, aos módulos que de fato apresentavam o problema, preservando o comportamento original no restante do sistema.

O segundo ajuste diz respeito ao esforço de raciocínio. O parâmetro correspondente na configuração do GraphRAG não produz efeito quando a chamada passa pela biblioteca intermediária litellm até o OpenRouter, pois é descartado sem aviso. Por isso, foi necessário injetar o parâmetro diretamente no corpo de cada requisição de geração, no formato reconhecido pelo OpenRouter, para que o esforço baixo fosse efetivamente aplicado.

O terceiro ajuste foi a paralelização das consultas, uma vez que a execução original, sequencial, era inviável em escala com um modelo de raciocínio lento. A concorrência foi ajustada empiricamente por modo de busca, a partir da medição do uso de memória por processo, já que o limite efetivo era a memória da máquina local, e não o limite de requisições da API. Adotaram-se seis processos simultâneos no modo \textit{basic}, três no modo \textit{local}, que consome cerca de 1,3~GB por processo, quase o dobro do \textit{basic}, e dois nos modos \textit{global} e \textit{drift}.

\section{Extração da resposta}\label{sec:met-parser}

Embora delimitada por marcas, a resposta de cada consulta é texto livre e precisa ser convertida em um conjunto de classes positivas. Essa conversão foi feita por um \textit{parser} determinístico, aplicado de forma idêntica a todos os sistemas e em todas as análises, a fim de garantir a comparabilidade dos resultados. O \textit{parser} percorre o trecho delimitado linha a linha em busca dos oito nomes canônicos e de variações mapeadas, com correspondência restrita a palavras inteiras. Uma linha só é contada como predição positiva se não contiver nenhum marcador de negação de uma lista auditada e corrigida ao longo do trabalho, que inclui, entre outros, ``ausente'', ``descartado'', ``negativo'', ``não confirmado'', ``nega'' e ``NR''.

Uma exceção a essa regra foi introduzida após uma auditoria de falsos negativos. Menções a carga viral negativa ou indetectável do vírus da hepatite B (HBV), registrada como PCR-HBV ou HBV-DNA, não são tratadas como negação da doença, pois indicam tratamento supressor bem-sucedido em paciente já diagnosticado com hepatite B crônica: a terapia da hepatite B é supressora, e não curativa, e o paciente permanece com a doença mesmo com carga viral indetectável. A exceção não se estende à hepatite C, cujo tratamento com antivirais de ação direta é frequentemente curativo, de modo que um PCR negativo para o vírus da hepatite C (HCV) após o tratamento costuma indicar doença de fato resolvida. Um teste realizado antes de restringir a correção à hepatite B confirmou que aplicá-la também à hepatite C produzia falsos positivos.

\section{Métricas e análise estatística}\label{sec:met-metricas}

Para cada sistema e cada classe $c$, os pares formados por paciente e doença foram classificados em verdadeiros positivos ($\mathrm{VP}_c$), falsos positivos ($\mathrm{FP}_c$), falsos negativos ($\mathrm{FN}_c$) e verdadeiros negativos ($\mathrm{VN}_c$), por comparação entre o conjunto previsto e os rótulos de referência. A partir dessas contagens, calcularam-se a precisão, a revocação (\textit{recall}) e o F1 de cada classe, conforme as definições da Seção~\ref{sec:avaliacao}. Essas métricas foram agregadas de duas formas, sendo $C$ o conjunto das oito classes. A agregação macro calcula a média simples entre as classes, atribuindo a todas o mesmo peso independentemente da prevalência, conforme a Equação~\eqref{eq:f1-macro}, e de forma análoga para a precisão e a revocação:
\begin{equation}
\mathrm{F1}_{\text{macro}} = \frac{1}{|C|}\sum_{c \in C} \mathrm{F1}_c .
\label{eq:f1-macro}
\end{equation}

A agregação micro, por sua vez, soma as contagens de todas as classes antes da divisão, ponderando implicitamente cada classe por sua frequência, conforme a Equação~\eqref{eq:pr-micro}; o $\mathrm{F1}_{\text{micro}}$ é a média harmônica entre $P_{\text{micro}}$ e $R_{\text{micro}}$:
\begin{equation}
P_{\text{micro}} = \frac{\sum_{c \in C} \mathrm{VP}_c}{\sum_{c \in C} \left(\mathrm{VP}_c + \mathrm{FP}_c\right)},
\qquad
R_{\text{micro}} = \frac{\sum_{c \in C} \mathrm{VP}_c}{\sum_{c \in C} \left(\mathrm{VP}_c + \mathrm{FN}_c\right)} .
\label{eq:pr-micro}
\end{equation}

Além das métricas por classe, adotou-se a acurácia de correspondência exata (\textit{exact match}, também denominada \textit{subset accuracy}), definida na Equação~\eqref{eq:em} como a proporção de pacientes cujo conjunto previsto $\hat{Y}_i$ é idêntico ao conjunto de referência $Y_i$, em que $N$ é o número de pacientes e $\mathbf{1}[\cdot]$ é a função indicadora. Trata-se da métrica mais rigorosa do trabalho, pois um único erro em qualquer das oito classes invalida o acerto do paciente inteiro, e é também a variável comparada no teste de McNemar (Seção~\ref{sec:met-mcnemar}).
\begin{equation}
\mathrm{EM} = \frac{1}{N}\sum_{i=1}^{N} \mathbf{1}\left[\hat{Y}_i = Y_i\right]
\label{eq:em}
\end{equation}

A acurácia por rótulo, definida na Equação~\eqref{eq:ar}, considera cada combinação de paciente e doença como um acerto ou erro binário independente, em que $\hat{y}_{ic}$ e $y_{ic}$ indicam, respectivamente, a previsão e o rótulo de referência da classe $c$ para o paciente $i$. Essa métrica é reportada com uma ressalva explícita: como a maior parte das combinações é negativa, já que cada paciente tem em média 1,49 das oito doenças, ela é inflacionada por verdadeiros negativos triviais e não deve ser lida isoladamente como medida de qualidade, em linha com a discussão sobre acurácia sob desbalanceamento da Seção~\ref{sec:avaliacao}. Para complementar as métricas agregadas, foram construídas matrizes de confusão $2 \times 2$ por doença e por sistema, visualizadas com a biblioteca matplotlib, que permitem identificar rapidamente qual classe concentra os erros de cada configuração.
\begin{equation}
\mathrm{AR} = \frac{1}{N\,|C|}\sum_{i=1}^{N}\sum_{c \in C} \mathbf{1}\left[\hat{y}_{ic} = y_{ic}\right]
\label{eq:ar}
\end{equation}

\subsection{Investigação do comportamento do modo \texorpdfstring{\textit{drift}}{drift}}\label{sec:met-hipoteses}

A comparação entre os sistemas revelou um padrão sistemático no comportamento do modo \textit{drift}, relatado no capítulo de resultados, que motivou uma investigação dedicada ao mecanismo subjacente. Trata-se, portanto, de uma análise conduzida \textit{a posteriori}, cujo desenho se descreve aqui. Foram comparados dois grupos de pares formados por prontuário e doença presente nos rótulos de referência. O grupo A reúne os casos em que o modo \textit{drift} identifica corretamente a doença e a maioria dos sistemas \textit{basic} e \textit{local} não a identifica; o grupo B reúne todos os demais casos positivos e funciona como controle.
% TODO: especificar se o acerto do drift considera S7, S8 ou ambos e qual o critério de "maioria" entre S1 a S4.

Quatro hipóteses foram formuladas, cada uma associada a um teste estatístico formal, implementado com o módulo de estatística da biblioteca SciPy:
\begin{description}
    \item[H1:] os prontuários do grupo A são mais longos, em número de caracteres, que os do grupo B; teste U de Mann-Whitney;
    \item[H2:] o termo literal da doença aparece em menor proporção no texto dos prontuários do grupo A; teste qui-quadrado de comparação de proporções;
    \item[H3:] os prontuários do grupo A registram mais atendimentos; teste U de Mann-Whitney;
    \item[H4:] a entidade canônica da doença deixou de ser extraída na indexação, para o paciente correspondente, em maior proporção no grupo A; teste qui-quadrado.
\end{description}
% TODO: citar Mann e Whitney (1947) e o teste qui-quadrado de Pearson; não há entradas no .bib.

A escolha dos testes segue a natureza de cada variável. Comprimento do texto e número de atendimentos são variáveis quantitativas sem pressuposto de normalidade, razão pela qual se empregou o teste U de Mann-Whitney \citep{mann_whitney1947}, não paramétrico, que dispensa esse pressuposto. A presença do termo literal e a ausência da entidade canônica são variáveis binárias, comparadas entre os grupos pelo teste qui-quadrado de Pearson \citep{pearson1900}. Adotou-se, também aqui, nível de significância de 5\%. A hipótese H4 originou-se da análise detalhada do subgrafo de um paciente específico, descrita na seção seguinte, e foi submetida ao teste justamente para verificar se o mecanismo observado nesse caso se generaliza para a amostra.

\section{Análise qualitativa de erros e visualização do grafo}\label{sec:met-qualitativa}

Como apoio à investigação qualitativa dos erros, foram gerados subgrafos de entidades e relações por paciente, desenhados com as bibliotecas NetworkX e matplotlib em leiaute de força dirigida de Fruchterman-Reingold e coloridos por comunidade.
% TODO: citar Fruchterman e Reingold (1991); não há entrada no .bib.
Essas visualizações serviram para ilustrar mecanismos de erro específicos, como o de uma entidade sinônima de uma doença alocada em comunidade diferente daquela que contém a entidade correspondente ao nome canônico, no mesmo paciente, situação compatível com o caráter puramente topológico do agrupamento descrito na Seção~\ref{sec:met-grafo}.

O procedimento foi depois generalizado para todos os 537 prontuários, com uma imagem por paciente, e complementado por uma comparação visual dos seis maiores subgrafos da amostra e por uma análise estrutural agregada da distribuição de entidades, relações e comunidades por prontuário. Essa análise tem caráter exploratório: orienta a formulação de hipóteses, como a H4, que só são tomadas como gerais após o teste estatístico correspondente.

\section{Limitações metodológicas}\label{sec:met-limitacoes}

O desenho descrito apresenta limitações que condicionam a interpretação dos resultados. A primeira decorre do modelo de linguagem: por ser de raciocínio híbrido, o \texttt{ling-3.0-flash} introduz alta variância de latência e de custo, o que compromete a reprodutibilidade do tempo de execução e, em pelo menos um caso, inviabilizou a conclusão, dentro de prazo razoável, de uma validação experimental planejada, relatada no capítulo de resultados. A segunda refere-se à PBE, cujos 2 casos positivos são insuficientes para qualquer métrica estatisticamente confiável sobre essa classe.

A terceira limitação diz respeito aos rótulos de referência, que podem conter discordâncias pontuais não resolvidas. Durante a correção do \textit{parser} de negação, identificou-se ao menos um caso em que a predição do modelo parecia clinicamente mais correta que o próprio rótulo de referência. A quarta refere-se à amostra, enriquecida por triagem textual: a prevalência das classes nela observada não deve ser lida como estimativa epidemiológica da população do ambulatório, e o grupo sem doença não equivale a um controle aleatório, como discutido na Seção~\ref{sec:met-amostra}.

A quinta limitação é a ausência de resultados confiáveis para os sistemas do modo \textit{global} (S5 e S6). O mecanismo de seleção dinâmica de comunidades pode gerar um volume desproporcional de chamadas ao modelo, de até cerca de 547 por consulta no pior caso, quando o nível inicial da hierarquia não contém comunidades relevantes e a busca recorre integralmente ao nível seguinte. Tentativas de ajustar os parâmetros desse mecanismo não resolveram o problema sem custo relevante de qualidade. Por essa razão, esses sistemas não foram concluídos de forma confiável com o modelo adotado, e a comparação efetivamente realizada contempla os modos \textit{basic}, \textit{local} e \textit{drift}.

A sexta limitação diz respeito à própria partição em comunidades. Em grafos de conhecimento esparsos, a otimização da modularidade admite um número muito grande de partições quase ótimas, de modo que sementes aleatórias distintas ou pequenas perturbações nas arestas podem produzir comunidades estruturalmente diferentes, o que torna as comunidades do Leiden intrinsecamente não reprodutíveis nesse tipo de grafo \citep{hossain2026}. Como o índice foi construído uma única vez, os resultados por nível de comunidade refletem uma realização particular dessa partição, e a estabilidade das comunidades entre execuções da indexação não foi avaliada.

% =====================================================================
% TODO: PENDÊNCIAS DESTE CAPÍTULO
% 1. Versão do prompt usada em S7 e S8 (drift) na rodada final.
% 2. Número de anotadores da "segunda opinião", formação e medida de concordância.
% 3. Parecer do CEP/CAAE e procedimento de anonimização.
% 4. S5 e S6 (global): confirmar uso de seleção dinâmica e como o nível 1 ou 2 se aplica.
% 5. N=531 no McNemar do drift contra N=537 nos demais: explicar.
% 6. [RESOLVIDO 2026-11-01] Instituição de origem dos dados: Hospital Universitário
%    da Universidade Federal de Alagoas (HU-UFAL), confirmado pelo autor.
% 7. Índice: corpus completo ou amostra; LLM e modelo de embeddings da indexação; tamanho e sobreposição dos chunks.
% 8. Modelo de LLM usado na triagem (etapa 2 do funil); parâmetros de decodificação na rodada final.
% 9. Modo basic: conferir com a descrição da Seção~\ref{sec:rag} (recuperação vetorial) e se o nível de comunidade é inócuo nele.
% 10. Critério do grupo A (S7, S8 ou ambos; limiar de "maioria" entre S1 a S4).
% 11. [PARCIALMENTE RESOLVIDO 2026-10-19] McNemar, Mann-Whitney e qui-quadrado de
%     Pearson já têm entrada no .bib (mcnemar1947, mann_whitney1947, pearson1900) e
%     são citados no texto. Ainda faltam: Fruchterman e Reingold (1991), LGPD.
% 12. (2026-10-16) Revisar esta correção de modelo (ling-2.6-flash na rodada final,
%     ling-3.0-flash só na tentativa de migração) contra qualquer outra citação
%     cruzada ao modelo em outros capítulos (ex.: resumo, introdução) que possa
%     ainda assumir ling-3.0-flash como modelo da rodada final.
%
% TODO: ITENS DA FONTE A MIGRAR PARA O CAPÍTULO DE RESULTADOS E DISCUSSÃO
% - Resultado do McNemar por modo (basic N=537, 45 contra 40, p=0,6646; local N=537,
%   44 contra 47, p=0,8341; drift N=531, 44 contra 61, p=0,1180) e sua interpretação
%   como justificativa a posteriori para manter os dois níveis.
% - S8 superior a S7 no F1 macro (0,798 contra 0,745).
% - Padrão agregado do drift: revocação de cerca de 95% contra 60 a 67% em basic/local;
%   precisão de cerca de 65 a 76% contra 78 a 86%.
% - Resultados de H1 a H4 (medianas de 6.160 contra 3.230 caracteres; termo literal em
%   69,4% contra 89,4%; mediana de 2 atendimentos em ambos, p=0,21; entidade ausente em
%   46,6% contra 45,7%, p=0,89) e o estudo de caso individual que motivou H4.
% - Três grupos de doenças segundo o compromisso entre revocação e precisão do drift
%   (marcador binário claro; inferência indireta; exceção da HDA).
% - Tentativa inconclusiva de validação da versão 3 do prompt para HDA no drift
%   (4 pacientes, teto de 45 minutos atingido).
% =====================================================================
